\documentclass[11pt,oneside]{book}
\usepackage[
    paperwidth=6in,
    paperheight=9in,
    top=0.75in,
    bottom=0.75in,
    left=0.6in,
    right=0.65in,
    twoside=false,
]{geometry}
\input{../pri-end/T0_preamble_De.tex}
% T0_preamble_patches.tex — LOKALER PLATZHALTER
% Im Repo durch die offizielle Version ersetzen.

\usepackage{pdfpages}
% Kapitelzählung: arabisch statt Buchstaben (Preamble-Override)
\makeatletter
\renewcommand{\thechapter}{\arabic{chapter}}
\makeatother

% Statusmarker als Marginalien
\newcommand{\status}[1]{\penalty10000\,{\tiny\textcolor{gray!70}{[#1]}}}

% Laierspur-Umgebung
\newenvironment{laierspur}{%
  \begin{tcolorbox}[
    colback=blue!3!white, colframe=blue!50!black,
    title={\small\textit{Für alle Leser}},
    fonttitle=\small\itshape,
    left=5pt, right=5pt, top=3pt, bottom=3pt, breakable
  ]\small\itshape}{%
  \end{tcolorbox}}

\setcounter{tocdepth}{0}

\begin{document}
\frontmatter
\pagestyle{fancy}
\makeatletter\let\ps@plain\ps@fancy\makeatother

\begin{titlepage}
  \centering\vspace*{1.5cm}
  {\Huge\bfseries\color{blue!70!black}
    Jenseits der Matrix}\\[0.6cm]
  {\Large\bfseries\color{blue!40!black}
    Die algebraische Struktur der Wirklichkeit}\\[1cm]
  {\LARGE\itshape Fundamentale Fraktale Geometrische\\[0.2cm]
    Feldtheorie (FFGFT)}\\[0.5cm]
  {\large\itshape Von $T^4/\mathbb{Z}_3$ zu den Teilchen des Standardmodells}\\[2cm]
  {\large Johann Pascher}\\[0.3cm]
  {\normalsize ORCID: 0009-0000-6518-4064}\\[0.3cm]
  {\normalsize\url{https://github.com/jpascher/T0-Time-Mass-Duality}}\\[1.5cm]
  {\normalsize September 2026}
  \vfill{\small Version 1.0 — Populärwissenschaftliche Ausgabe}
\end{titlepage}

\newpage\thispagestyle{empty}\vspace*{\fill}
{\small\noindent\textbf{Bibliografische Information}\\[0.4em]
Johann Pascher: \textit{Jenseits der Matrix. Die algebraische Struktur der Wirklichkeit. FFGFT.}\\
Erste Ausgabe, September 2026.\\[0.6em]
\noindent Zenodo: DOI 10.5281/zenodo.17390357\\
Lizenz: Creative Commons BY 4.0}\vspace*{\fill}

\tableofcontents
\mainmatter

\part{Die Frage}

\chapter{Leben wir in einer Matrix?}

Die Frage ist alt. Platon stellte sie mit dem Gleichnis der Höhle. Descartes als radikalen Zweifel. Und 2003 gab Nick Bostrom ihr eine moderne Form: Wenn fortgeschrittene Zivilisationen Simulationen laufen lassen können, leben wir mit überwältigender Wahrscheinlichkeit in einer solchen Simulation.

\begin{laierspur}
Stell dir vor, du spielst ein Videospiel, das so perfekt ist, dass du von innen nicht merken kannst, ob es ein Spiel ist. Die Physik wäre konsistent, die Mathematik stimmte --- du hättest keine Möglichkeit zu unterscheiden. Das ist Bostroms Argument.
\end{laierspur}

Es ist eine ernste Frage. Und sie ist falsch gestellt. \enquote{Simulation} ist die falsche Kategorie --- sie setzt einen Simulanten voraus, eine äußere Instanz.

Was wenn das Universum keine Simulation \textit{ist}, sondern eine mathematische Struktur \textit{hat} --- tief, zwingend, ohne äußere Instanz? Was wenn die Frage nicht lautet: \textit{wer rechnet uns}, sondern: \textit{was ist die Struktur, die die Natur nicht anders kann?}

Was diese Schrift vorlegt, ist eine solche Struktur: ein kompakter vierdimensionaler Torus $T^4$ mit einer dreifachen Symmetrie $\mathbb{Z}_3$, aus der algebraisch folgt, was folgen muss. Keine freien Parameter. Kein Postulat über Teilchen.

\chapter{Warum nicht zehn?}

Das Standardmodell hat neunzehn freie Parameter --- Teilchenmassen, Kopplungsstärken, Mischungswinkel. Es nimmt sie hin. Darunter sticht eine Tatsache heraus: Es gibt drei Generationen von Fermionen, neun fundamentale Familien.

\begin{laierspur}
Stell dir vor, du findest eine Uhr. Sie läuft präzise --- aber du weißt nicht, warum sie so gebaut ist. Warum zwölf Stunden, warum diese Zahnräder? Das Standardmodell ist diese Uhr. FFGFT fragt: nach welchem Prinzip wurde sie gebaut?
\end{laierspur}

Das LEP-Experiment fand keine vierte leichte Neutrinogeneration --- experimentell, durch Messung. Aber \enquote{nicht gefunden} ist nicht \enquote{strukturell unmöglich}.

FFGFT gibt eine algebraische Antwort: Der Raum $T^4/\mathbb{Z}_3$ hat genau neun Fixpunkte. Nicht acht. Nicht zehn. Neun.\status{[B] Dok.~282}

\textbf{Eine zehnte Familie ist algebraisch verboten.}

\begin{laierspur}
Wenn du ein Dreieck zeichnest, hat es genau drei Ecken --- nicht weil du aufgehört hast, sondern weil die Struktur es so bestimmt. Der $\mathbb{Z}_3$-Zirkulant hat drei Eigenwerte. Eine vierte Generation würde nicht eine vierte Ecke hinzufügen --- sie würde die Struktur zerstören.
\end{laierspur}

Eine Theorie die sagt \enquote{es gibt neun} ist beeindruckend. Eine Theorie die sagt \enquote{es können keine anderen sein} ist eine andere Klasse von Aussage.

\part{Die Geometrie}

\chapter{Was Mathematik der Natur schuldet}

Drei philosophische Positionen stehen zur Wahl: Instrumentalismus (Mathematik als Werkzeug), Strukturalismus (Tegmark: das Universum \textit{ist} Mathematik), und geometrischer Konstitutionalismus: eine geometrische Grundstruktur erzwingt algebraische Konsequenzen, die mit Beobachtung übereinstimmen --- ohne Wesensaussage.

FFGFT nimmt die dritte Position. Ihr Statussystem unterscheidet:

\begin{description}[leftmargin=3em]
  \item[\textbf{[SETZUNG]}] Axiom --- deklariert, nicht hergeleitet.
  \item[\textbf{[B]}] Algebraisch bewiesen, Prüfskript verifiziert.
  \item[\textbf{[K]}] Aus $\xi$ hergeleitet, numerisch geprüft.
  \item[\textbf{[S]}] Skizziert --- plausibel, Ableitung noch offen.
  \item[\textbf{[X]}] Widerlegt durch Beweis oder Prüfskript.
\end{description}

\begin{laierspur}
Ein guter Tischler sagt, welche Verbindung verleimt ist und welche nur gesteckt. Das Statussystem macht FFGFT auditierbar --- nicht weil die Theorie schwach ist, sondern weil sie ehrlich ist.
\end{laierspur}

\chapter{Der Torus und $\tilde{T} \cdot m = 1$}

Ein Torus entsteht durch Identifikation: rechter Rand = linker Rand, oberer = unterer. Lokal flach, global geschlossen.

\begin{laierspur}
Ein altes Computerspiel, bei dem die Figur rechts verschwindet und links wiederkommt. Lokal normal --- global verbunden. Das ist $T^4/\mathbb{Z}_3$.
\end{laierspur}

Aus der Kompaktifizierung eines der vier Kreise folgt die Grundrelation:\status{[K] Dok.~077}
\[
  \tilde{T} \cdot m = 1
\]
Masse ist die inverse Zeitskala einer Mode. Schwer = kurze intrinsische Zeit. In natürlichen Einheiten ist $E = mc^2 = E = m$ --- der Faktor $c^2$ ist ein Einheitenkonversionsfaktor, kein Naturgesetz.

\chapter{Die Dreiteilung: $\mathbb{Z}_3$ und die neun Fixpunkte}

$\mathbb{Z}_3$ hat drei Elemente: Identität, Rotation um $120°$, Rotation um $240°$. Die Wirkung von $\mathbb{Z}_3$ auf $T^4$ erzeugt \textbf{neun Fixpunkte}.\status{[B] Dok.~282}

Der hermitesche $\mathbb{Z}_3$-Zirkulant auf $\mathbb{C}^3$ hat exakt drei Eigenwerte $\lambda_k = a + b\omega^k + c\omega^{2k}$ ($\omega = e^{2\pi i/3}$) --- drei Generationen, nicht mehr.

Die Frobenius-Trennung in $GF(27)^*$:\status{[B] Dok.~339}
\begin{itemize}
  \item Fixpunkte $\{+1,-1\}$ = massiver Sektor (Leptonen, Quarks)
  \item Acht Dreier-Orbits = acht Gluonen
  \item $\{f_1, f_2\}$ = Majorana ($\nu_1, \nu_2$); $\{f_3, f_4\}$ = Dirac
  \item $\nu_3$ ist Dirac --- \textbf{falsifizierbare Vorhersage}\status{[B] Dok.~343}
\end{itemize}

\part{Die Algebra}

\chapter{Endliche Körper}

$GF(q)$: endlicher Körper mit $q = p^n$ Elementen. Die physikalische Hierarchie:

\begin{center}
\begin{tabular}{lll}
\toprule
Körper & Ordnung & Physikalische Rolle \\
\midrule
$GF(3)$ & 3 & Farbladung \\
$GF(9)^*$ & 8 & Acht Gluonen \\
$GF(27)^*$ & 26 & Vollständige Fermionstruktur \\
\bottomrule
\end{tabular}
\end{center}

Das Legendre-Symbol $(k/13) \in \{+1,-1\}$ erzwingt Ladungsquantisierung $Q \in \{0, \pm\tfrac{1}{3}, \pm\tfrac{2}{3}, \pm1\}$.\status{[B] Dok.~346}

\chapter{Das Galois-Raster der Teilchenmassen}

Massenformel: $m_i = r_i \cdot \xi^{p_i} \cdot v \cdot K_\text{frak}$\status{[K] Dok.~006}

Wicklungszahlen der Leptonen:
\begin{center}
\begin{tabular}{lcccc}
\toprule
Teilchen & $n_\theta$ & $n_\phi$ & $r_i$ & $p_i$ \\
\midrule
$e$ & 3 & 2 & $4/3$ & $3/2$ \\
$\mu$ & 5 & 4 & $16/5$ & 1 \\
$\tau$ & 9 & 5 & $25/9$ & $2/3$ \\
\bottomrule
\end{tabular}
\end{center}

Die Schlüsselidentität:\status{[K] Dok.~338}
\[
  \frac{(r_\mu/r_e)^2}{\xi} = |GF(9)^*|^2 \cdot 5^2 \cdot |GF(27)| = 64 \cdot 25 \cdot 27 = 43200
\]
Daraus folgt $\xi = 4/30000$ --- algebraisch erzwungen, kein freier Parameter.

\chapter{Drei Generationen: Verbot der vierten}

Der $\mathbb{Z}_3$-Zirkulant ist eine $3\times3$-Matrix --- drei Eigenwerte, nicht vier. Strukturelles Verbot der vierten Generation.\status{[B] Dok.~282}

Ladungsquantisierung und Farbe aus zwei Galois-Eigenschaften:\status{[B] Dok.~346}
\begin{center}
\begin{tabular}{lll}
\toprule
QR mod 13 & $N\bmod3$ & Teilchentyp \\
\midrule
QR (neutral) & 0 (Singlett) & Neutrino \\
NQR (geladen) & 0 (Singlett) & geladenes Lepton \\
NQR (geladen) & 1,2 (Triplett) & Quark \\
\midrule
QR + Triplett & --- & \textbf{verboten} \\
\bottomrule
\end{tabular}
\end{center}

\chapter{Gluonen, Farbe und Confinement}

Die starke Kernkraft wird durch acht Gluonen vermittelt. Warum acht?
Das Standardmodell setzt es als Eigenschaft der Eichgruppe $SU(3)$.
FFGFT leitet es her.

\section*{Gluonen als Galois-Objekte}

In $GF(27)^*$ wirkt der Frobenius $\phi: x\mapsto x^3$ und zerlegt die
26 Elemente in Fixpunkte $\{+1,-1\}$ (massiver Sektor) und
Dreier-Orbits. Die vier $\mathbb{Z}_{13}$-Orbits kombiniert mit
$\mathbb{Z}_2$-Parität ergeben:\status{[B] Dok.~339}
\[
  4 \text{ Orbits} \times 2 \text{ Paritäten} = \mathbf{8 \text{ Gluonen}}
\]
Die drei Elemente jedes Orbits sind die drei Farben
$\{\mathrm{R},\mathrm{G},\mathrm{B}\}$. Der Fixkörper $GF(3)^* = \{1,2\}
\cong\mathbb{Z}_2$ liefert das Photon.

\begin{center}
\small
\begin{tabular}{ll}
\toprule
Strukturelement & Physikalische Bedeutung \\
\midrule
Fixpunkte $\{+1,-1\}$ & Massiver Sektor (Fermionen) \\
4 Orbits $\times$ 2 Paritäten & Acht Gluonen \\
3 Elemente je Orbit & Drei Farben R, G, B \\
Fixkörper $GF(3)^*$ & Photon ($U(1)$-Sektor) \\
\bottomrule
\end{tabular}
\end{center}

\begin{laierspur}
Stell dir 26 Elemente vor, aufgeteilt in acht Dreiergruppen plus zwei
Fixpunkte. Die acht Gruppen sind die Gluonen, die drei Positionen die
Farben, die zwei Fixpunkte die massive Materie. Alles aus einer einzigen
Gruppenstruktur --- kein freier Parameter.
\end{laierspur}

\section*{Confinement als $\mathbb{Z}_3$-Trialitätsselektion}

Warum hat niemand je ein einzelnes Quark gesehen? In FFGFT ist die
Antwort algebraisch: Nur $\mathbb{Z}_3$-\textbf{invariante} Objekte
--- mit Gesamttrialität $N \bmod 3 = 0$ --- können den Torus verlassen
und beobachtet werden.\status{[B] Dok.~325}

\begin{itemize}
  \item Einzelnes Quark: $N \bmod 3 = 1$ oder $2$ --- \textbf{nicht beobachtbar}
  \item Drei Quarks (Baryon): $1+1+1 \equiv 0$ --- \textbf{beobachtbar}
  \item Quark+Antiquark (Meson): $1+2 \equiv 0$ --- \textbf{beobachtbar}
\end{itemize}

\begin{laierspur}
Eine Türsteher-Regel: Nur Gruppen, deren Summe ein Vielfaches von drei
ist, kommen rein. Ein Quark mit Ticket ``1'' kommt nicht rein. Drei
Quarks mit ``1'', ``1'', ``1'' --- Summe 3 --- kommen rein. Das ist
$\mathbb{Z}_3$-Trialität. Das ist Confinement.
\end{laierspur}

\section*{Das Proton}

Ein Proton (uud) trägt die Farben $\{\mathrm{R},\mathrm{G},\mathrm{B}\}$
--- $\mathbb{Z}_3$-invariant. Die Ladung $+1 = 2\times(+2/3) + 1\times(-1/3)$
folgt direkt aus der Ladungsquantisierung via Legendre-Symbol.\status{[B] Dok.~346}

Die Protonmasse 938 MeV entsteht zum größten Teil aus dem Gluonenfeld
(Gitter-QCD), nicht aus den Quarkmassen. Die Galois-Herleitung der
genauen Masse ist offen.\status{[S]}

\chapter{Antiteilchen, Spin und Chiralität: Zustände, keine Verdopplung}

Wenn man die Teilchen des Standardmodells zählt, stößt man schnell auf
eine Verwirrung. Es gibt zwölf Fermionen --- aber dann gibt es auch
Antiteilchen: das Positron zum Elektron, das Antiquark zum Quark.
Macht das 24? Und dann noch Spin-up und Spin-down für jedes Teilchen.
Macht das 48?

Die Antwort ist: nein. Und die Erklärung liegt in der Algebra.

\section*{Was ein Antiteilchen ist}

Ein Antiteilchen ist nicht ein anderes Teilchen. Es ist dasselbe
algebraische Objekt --- mit umgekehrter Ladung. Das Positron hat dieselbe
Masse wie das Elektron, denselben Spin, dieselbe Wechselwirkungsstärke.
Nur die elektrische Ladung ist entgegengesetzt: $+1$ statt $-1$.

In $GF(27)^*$ entspricht dieser Unterschied der $\mathbb{Z}_2$-Parität ---
dem Faktor $\mathbb{Z}_2$ in der Zerlegung $\mathbb{Z}_{26} \cong
\mathbb{Z}_2 \times \mathbb{Z}_{13}$. Teilchen und Antiteilchen sind
die zwei $\mathbb{Z}_2$-Elemente desselben Orbits. \status{[B] Dok.~339}

Es sind keine zwei getrennten Objekte. Es ist ein algebraisches Objekt
mit zwei Zuständen --- wie eine Münze mit zwei Seiten, nicht wie zwei
verschiedene Münzen.

\begin{laierspur}
Stell dir eine Uhr vor. Die Zeiger können im oder gegen den Uhrzeigersinn
drehen. Das ist kein zweites Uhrwerk --- es ist dasselbe Uhrwerk in zwei
möglichen Zuständen. Teilchen und Antiteilchen sind dasselbe algebraische
Objekt, einmal mit Ladung $+q$, einmal mit $-q$. Die $\mathbb{Z}_2$-Parität
in $GF(27)^*$ ist genau dieser Umkehrzustand.
\end{laierspur}

\section*{Was Spin ist}

Spin ist noch fundamentaler. Er ist kein rotierender Ball --- das war
das historische Bild, das schon 1925 scheiterte, weil es unphysikalische
Rotationsgeschwindigkeiten ergeben hätte. Spin ist eine algebraische
Eigenschaft, die aus der Struktur des Raumes folgt.

In der Clifford-Algebra Cl(6) --- der algebraischen Beschreibung des
sechsdimensionalen internen Raumes der FFGFT --- hat jeder Zustand einen
Grad $r$, der bestimmt, wie er unter Rotationen transformiert. Spin-1/2-
Teilchen entsprechen Grad-1-Objekten (Vektoren) im Clifford-Raum.

Das Besondere: Ein Spin-1/2-Teilchen braucht eine \textbf{Drehung um
720°}, um in sich selbst überzugehen --- nicht 360°. Das ist kein
Paradox. Es folgt direkt aus der algebraischen Struktur:

\begin{center}
$\gamma_5 \cdot X = (-1)^r \cdot X$
\end{center}

Für ungerade Grade $r$: $\gamma_5 X = -X$ --- das ist
\textbf{linkshändig}. Für gerade Grade: $+X$ ---
\textbf{rechtshändig}. \status{[B] Dok.~349}

Spin-up und Spin-down sind ebenfalls keine zwei verschiedenen Teilchen.
Sie sind zwei Projektionen desselben algebraischen Objekts auf verschiedene
Raumrichtungen.

\begin{laierspur}
Ein Schatten an der Wand kann lang oder kurz sein, je nachdem wie das
Licht fällt. Der Gegenstand ist derselbe --- der Schatten ist eine
Projektion. Spin-up und Spin-down sind Projektionen desselben Spinors
auf verschiedene Messrichtungen. Keine zwei Objekte --- ein Objekt,
zwei Messzustände.
\end{laierspur}

\section*{Was Chiralität ist}

Die schwache Kernkraft verhält sich seltsam: Sie wirkt nur auf
linkshändige Teilchen und auf rechtshändige Antiteilchen. Das ist eine
der merkwürdigsten Asymmetrien der Natur --- die sogenannte
\textbf{Paritätsverletzung}, entdeckt 1956.

In FFGFT folgt diese Asymmetrie aus der Cl(6)-Gradparität --- aus
demselben $\gamma_5$-Operator, der Spin definiert. Der
\textbf{Su-Sektor} (gerade Cl(6)-Grade) ist rechtshändig. Der
\textbf{Sd-Sektor} (ungerade Cl(6)-Grade) ist linkshändig. Und die
schwache Kraft $SU(2)_L$ koppelt nur an den Sd-Sektor --- algebraisch
erzwungen, nicht gesetzt. \status{[B] Dok.~349}

Das rechtshändige Neutrino $\nu_R$ ist dabei strukturell absent:
Im Sd-Sektor gibt es keinen Platz dafür. Das ist konsistent mit der
Vorhersage aus Kapitel 10 (Neutrinos): $\nu_3$ ist Dirac, aber
$\nu_R$ als eigenständiger Freiheitsgrad existiert nicht in
$GF(27)^*$. \status{[B]/[S] Dok.~349}

\begin{laierspur}
Stell dir eine Schere vor. Sie hat eine linke und eine rechte Seite ---
aber es ist eine Schere, nicht zwei. Die schwache Kraft greift nur
auf der linken Seite an. In FFGFT folgt diese Einseitigkeit aus der
algebraischen Gradparität: gerade Grade = rechts, ungerade = links.
Die Natur ist nicht symmetrisch zwischen links und rechts ---
und die Algebra erklärt warum.
\end{laierspur}

\section*{Die vollständige Zählung}

Was sind dann die „12 Teilchen" wirklich?

\begin{center}
\small
\begin{tabular}{lll}
\toprule
Eigenschaft & Algebraische Herkunft & Anzahl der Zustände \\
\midrule
Teilchenfamilien & 9 Fixpunkte von $T^4/\mathbb{Z}_3$ & 9 \\
Farbe & $N \bmod 3$ in $GF(3)$ & $\times 3$ für Quarks \\
Teilchen/Antiteilchen & $\mathbb{Z}_2$-Parität in $GF(27)^*$ & $\times 2$ \\
Chiralität & Cl(6)-Gradparität & $\times 2$ \\
Spin-Projektion & Messrichtung & 2 Zustände \\
\bottomrule
\end{tabular}
\end{center}

Nichts davon ist eine eigenständige Verdopplung. Alles sind Zustände
desselben algebraischen Objekts --- unterschiedliche Projektionen,
Paritäten, Grade innerhalb einer einzigen Struktur.

Die „12 Fermionen" des Standardmodells sind in FFGFT: 3 Leptonen
(QR-Sektor, farblos) + 3×3 Quarks (NQR-Sektor, 3 Farben) = 12
\textit{Familien-Farb-Kombinationen}. Antiteilchen, Spin und Chiralität
sind keine zusätzlichen Objekte --- sie sind algebraische Zustände
innerhalb derselben Galois-Struktur.

\chapter{Das Higgs-Boson: Teilchen oder Vakuum-Artefakt?}

Im Juli 2012 meldete das CERN die Entdeckung des Higgs-Bosons --- eines
Teilchens mit Masse 125 GeV, das das Standardmodell seit 1964 vorhergesagt
hatte. Ein Triumph. Aber eine Frage blieb offen, die im Jubel unterging:

\textit{Was ist das Higgs-Boson wirklich?}

\section*{Was das Standardmodell sagt}

Im Standardmodell gibt es ein Higgs-Feld $\Phi$, das den gesamten Raum
durchdringt. Dieses Feld hat einen \textbf{Vakuumerwartungswert}
$\langle\Phi\rangle = v \approx 246$ GeV --- es ist im Vakuum nicht null.
Die Wechselwirkung der Teilchen mit diesem Hintergrundfeld erzeugt ihre
Masse: $m = y \cdot v$, wobei $y$ die Yukawa-Kopplung ist.

Das Higgs-Boson $H^0$ ist die Quantenschwingung dieses Feldes um seinen
Gleichgewichtswert --- die Anregung, die 2012 am LHC beobachtet wurde.

Warum hat das Vakuum diesen nichtverschwindenden Wert? Weil das
Higgs-Potential $V(\Phi) = \mu^2|\Phi|^2 + \lambda|\Phi|^4$ mit $\mu^2 < 0$
ein Minimum nicht bei $\Phi = 0$ hat, sondern bei $|\Phi| = v$.
Und warum ist $\mu^2 < 0$? Das Standardmodell antwortet: \textit{das
wird gesetzt. So ist es eben.}

\begin{laierspur}
Stell dir ein Tal vor. Das Standardmodell sagt: Das Universum liegt im
Tal, weil das Tal dort ist. Warum ist das Tal dort? Weil das Potential
diese Form hat. Warum hat es diese Form? Keine Antwort.
FFGFT fragt anders: Welche Geometrie erzwingt dieses Tal?
\end{laierspur}

\section*{Was FFGFT sagt}

In FFGFT ist $\tilde{T} \cdot m = 1$ keine Postulat über Felder ---
es ist eine Kompaktifizierungsrelation auf $T^4$. Und diese Relation
hat eine direkte Konsequenz: \textbf{Ein Vakuum mit $m = 0$ ist
algebraisch ausgeschlossen.} \status{[B] Dok.~354}

Wenn $m = 0$, dann $\tilde{T} \to \infty$ --- eine unendlich lange
intrinsische Zeitskala. Das ist physikalisch nicht realisierbar auf
einem kompakten Torus. Der Torus \textit{erzwingt} Masse im Vakuum.
Was das Standardmodell durch $\mu^2 < 0$ setzt, folgt in FFGFT aus
der Geometrie.

Das Higgs-Feld $\Phi$ und das Zeitfeld $\tilde{T}$ sind dabei
\textbf{invers zueinander}: \status{[B] Dok.~354}
\[
  \tilde{T}(x,t) = \frac{1}{m(x,t)} = \frac{1}{y \cdot (v + h(x,t))}
\]
Das Zeitfeld ist das inverse Higgs-Feld. Der Higgs-Mechanismus ist die
\textit{flache Projektion} der kompakten $T^4/Z_3$-Topologie in die
Sprache der Quantenfeldtheorie.

\begin{laierspur}
Eine Kugel, die auf einem Globus rollt, folgt einer Geodäte ---
der kürzesten Verbindung auf der gekrümmten Oberfläche. In der
flachen Landkarte sieht dieselbe Bahn gebogen aus, als ob eine
mysteriöse Kraft sie ablenkt. Diese ``Kraft'' ist kein echtes Prinzip ---
sie ist ein Artefakt der Projektion. Der Higgs-Mechanismus ist in FFGFT
ähnlich: eine echte Struktur (der Torus), projiziert auf eine flache
Beschreibung (das Higgs-Potential).
\end{laierspur}

\section*{Existiert das Higgs-Boson?}

Diese Frage hat in FFGFT eine präzise Antwort --- mit zwei Teilen.

\textbf{Teil 1: Das Higgs-Feld existiert.} Der nichtverschwindende
Vakuumerwartungswert $\langle\Phi\rangle \neq 0$ ist in FFGFT nicht
postuliert sondern erzwungen. Das Vakuum trägt Masse --- das ist
[B]. Die spontane Symmetriebrechung
$SU(2)_L \times U(1)_Y \to U(1)_\text{EM}$ folgt aus drei
Galois-Bausteinen: \status{[B] Dok.~355}
\begin{itemize}
  \item $\tilde{T}\cdot m=1$ erzwingt $\langle\Phi\rangle \neq 0$
  \item $Z_3$-Galois-Struktur erzwingt $Q_\text{vak}=0$ (neutrales Vakuum)
  \item $Q=0$ selektiert $U(1)_\text{EM}$ als Restgruppe
\end{itemize}

\textbf{Teil 2: Ob das Higgs-Boson $H^0$ eine eigenständige Anregung
ist oder ein Artefakt der flachen Beschreibung --- das ist offen.}
\status{[S] Dok.~354}

Das experimentell bei 125 GeV beobachtete Teilchen ist real [X].
Die Frage ist, was es \textit{ist}: Eine echte Anregung eines
fundamentalen Feldes? Oder die Darstellung einer Torus-Topologie in
der Sprache der flachen Quantenfeldtheorie --- nützlich als
Beschreibung, aber nicht fundamental?

\begin{laierspur}
Vergleich: Der Schall ist real --- wir hören ihn. Aber Schall ist
keine fundamentale Größe: er ist eine Druckwelle in Luft. In einer
Theorie ohne Luft gibt es keinen Schall. Das Higgs-Boson könnte
ähnlich sein: real als Beobachtung, aber nicht fundamental --- eine
Welle in einem Vakuumfeld, das selbst ein Projektionsartefakt ist.
FFGFT sagt: möglicherweise. Noch nicht bewiesen.
\end{laierspur}

Das ist die ehrlichste Aussage, die FFGFT machen kann: Der
Higgs-Mechanismus folgt aus der Geometrie [B]. Das Higgs-Boson
als Teilchen ist experimentell gesichert [X]. Ob es fundamental
oder Projektion ist --- das bleibt eine offene Frage [S].

\chapter{Neutrinos: Majorana oder Dirac?}

Majorana-Kriterium: $k \mapsto k^{-1} \bmod 13$. Orbit 3 = $\{4,10,12\}$ selbstinvers: Majorana. Orbit 4 = $\{7,8,11\}$ nicht selbstinvers: Dirac.\status{[B] Dok.~340}

Neutrinomassen:\status{[K] Dok.~340}
\begin{center}
\begin{tabular}{lll}
\toprule
Neutrino & Masse & Galois-Herkunft \\
\midrule
$\nu_1$ & $4{,}54$ meV & Orbit 1, $m_\nu$ \\
$\nu_2$ & $9{,}81$ meV & Orbit 3, $\sqrt{14/3}\,m_\nu$ \\
$\nu_3$ & $49{,}96$ meV & Orbit 4, $11\,m_\nu$ \\
\bottomrule
\end{tabular}
\end{center}

$\Delta m^2_\text{atm}$: Abw.~1,0\,\%; $\Delta m^2_\text{sol}$: Abw.~0,46\,\%.
Vorhersage: $m_{ee} \leq 14{,}4$ meV (testbar durch KamLAND-Zen).

\part{Die Konstanten}

\chapter{Ein Parameter, alle Konstanten}

\begin{center}
\begin{tabular}{llll}
\toprule
Konstante & FFGFT & PDG & Abw. \\
\midrule
$1/\alpha$ & $3700/27 = 137{,}037$ & $137{,}036$ & 7,6 ppm \\
$\sin^2\theta_W$ & $0{,}2305$ & $0{,}23122$ & $0{,}06\,\%$ \\
$m_\mu/m_e$ & $207{,}846$ & $206{,}768$ & $0{,}52\,\%$ \\
$\Delta m^2_\text{atm}$ & $2{,}476\times10^{-3}$ eV² & $2{,}500\times10^{-3}$ eV² & $1{,}0\,\%$ \\
$\Delta m^2_\text{sol}$ & $7{,}565\times10^{-5}$ eV² & $7{,}530\times10^{-5}$ eV² & $0{,}46\,\%$ \\
$\Sigma m_\nu$ & $64{,}3$ meV & $<120$ meV & konsist. \\
$m_\tau$ & $1776{,}78$ MeV & $1776{,}86\pm0{,}12$ MeV & $0{,}4\sigma$ \\
\bottomrule
\end{tabular}
\end{center}

\chapter{Leptonen und ihre Reste}

Reste $\varepsilon_i$ in Einheiten $\xi$: $+39\xi$, $+78\xi$, $+117\xi$ --- systematisch fallend.\status{[K] Dok.~352}

QED-Selbstenergie als Ursache: auf $112\sigma$ ausgeschlossen.\status{[K] Dok.~352}

Offene Brücke R113: Attribution offen.\status{[S] R113}

\chapter{Die Zeta-Funktion des Torus}

\[
  \zeta_{T^4/\mathbb{Z}_3}(s) = (1 - 3^{-s}) \cdot \zeta(s)
\]\status{[B] Dok.~343}

Neue Nullstellen bei $s = 2\pi ik/\ln 3$ auf Re$(s)=0$ --- Signatur der $\mathbb{Z}_3$-Symmetrie. Keine Resonanz mit Riemann-Nullstellen. FFGFT berührt die Riemann-Hypothese --- löst sie nicht.

\part{Repräsentation und Sein}

\chapter{Repräsentation, nicht Ontologie}

Die Galois-Klassen $\{f_1, f_2, f_3, f_4\}$ repräsentieren Teilchen --- sie sind keine Teilchen. Geometrie erzwingt Algebra. Algebra entspricht Beobachtung. Ontologie bleibt offen.

\chapter{Superposition ohne Zeit}

Kein Zeitoperator in der QM (Pauli 1933). Superposition ist zeitlos.

In FFGFT: $\tilde{T}_i = 1/m_i$ modengebunden. Modenübergreifende Superposition erzeugt inkompatible Zeitskalen.\status{[B] Dok.~334} Der Messakt verwirft die Windungszahl --- der Zeitpfeil hat einen geometrischen Grund.\status{[K] Dok.~333}

\chapter{Geometrie als tiefste Basis}

Konstitutionstyp: \textbf{Geometrie [SETZUNG] $\to$ Algebra [B] $\to$ Empirie [K]}.

Lorentz-Invarianz: lokal erhalten, global topologisch.\status{[K] Dok.~312}

Bell-Korrelationen als Holonomie auf $T^4$.\status{[B] Dok.~230}

\part{Rückblick und Ausblick}

\chapter{Was bewiesen, was skizziert, was offen ist}

\section*{Bewiesen [B]}
Neun Fixpunkte. Drei Generationen. Verbot der vierten. Acht Gluonen. Majorana-$\nu_{1,2}$. Dirac-$\nu_3$. Ladungsquantisierung. Gell-Mann-Nishijima. Verbot neutral+farbig. Kopplungsmatrix $f_3\times f_4$. Einträge $1/\sqrt{2}$ in $f_1\times f_3$. Alle Yukawa-Koeffizienten im Galois-Raster.

\section*{Offen [S]}
R113-Brücke (fallende Reste). $\tau_\mu$ direkt. Kosmischer Exponent P20. CMB-Peaks. Quarksektor vollständig.

\chapter{Die Matrix-Frage beantwortet}

\begin{description}
  \item[Simulation] Das Universum ist Artefakt --- kein Simulant in FFGFT.
  \item[Tegmark] Eine von unendlich vielen Strukturen --- FFGFT macht konkrete Vorhersagen.
  \item[Zuse/Wolfram] Diskrete Regeln --- geometrische Konsistenz fehlt dort.
  \item[FFGFT] Das Universum ist wie ein Beweis. Es kann nicht anders sein.
\end{description}

\begin{laierspur}
Simulation --- wie ein Computerspiel. Tegmark --- wie eine Zahl. Zuse --- wie ein Algorithmus. FFGFT --- wie ein Beweis. Es kann nicht anders sein.
\end{laierspur}

Falsifizierbare Vorhersagen:
\begin{itemize}
  \item $\nu_3$ ist Dirac (kein neutrinoloser Betazerfall für $\nu_3$)
  \item $\Sigma m_\nu \approx 64{,}3$ meV (KATRIN, Euclid)
  \item $m_\tau = 1776{,}78$ MeV (Belle II, BESIII)
  \item Keine vierte Quark-Generation (LHC)
\end{itemize}

Feynman schrieb $1/137$ an die Wand und fragte: Warum?

FFGFT schreibt: Weil $3700/27$.

\bigskip
\begin{center}\itshape
Das Universum ist keine Simulation. Es ist eine algebraische Struktur ---\\
und diese Struktur lässt keine anderen Teilchen zu als die, die wir kennen.\\
Das ist weniger als Ontologie. Und mehr, als die Physik bisher hatte.
\end{center}


\chapter*{Anhang A — Statussystem und Dokumentnummerierung}
\addcontentsline{toc}{chapter}{Anhang A — Statussystem und Dokumentnummerierung}

\section*{Statusmarker}
\begin{description}[leftmargin=4em]
  \item[\textbf{[SETZUNG]}] Axiom. In FFGFT: $T^4/\mathbb{Z}_3$.
  \item[\textbf{[B]}] Algebraisch bewiesen, Prüfskript verifiziert.
  \item[\textbf{[K]}] Aus $\xi$ hergeleitet, numerisch geprüft.
  \item[\textbf{[S]}] Skizziert --- plausibel, Ableitung offen.
  \item[\textbf{[X]}] Ausgeschlossen.
  \item[\textbf{[Q]}] Externe Primärquelle.
\end{description}

Im Buch erscheinen Marker als Marginalien (rechter Rand). Korrekturregister: Dok.~190.

\section*{Zugang}
\begin{itemize}
  \item GitHub: \url{https://github.com/jpascher/T0-Time-Mass-Duality}
  \item Zenodo: \href{https://doi.org/10.5281/zenodo.17390357}{DOI 10.5281/zenodo.17390357}
  \item ORCID: \href{https://orcid.org/0009-0000-6518-4064}{0009-0000-6518-4064}
\end{itemize}

\chapter*{Anhang B — Mathematische Grundlagen}
\addcontentsline{toc}{chapter}{Anhang B — Mathematische Grundlagen}

\section*{Galois-Felder}
$GF(p^n)$: endlicher Körper. $GF(3)=\{0,1,2\}$ mod 3. $GF(9)$: Erweiterung mit $i^2=2$. $GF(27)$: kubische Erweiterung.

\section*{Frobenius-Automorphismus}
$\phi: x \mapsto x^3$ auf $GF(27)$. Ordnung 3.
Fixpunkte $\{+1,-1\}$ = massiver Sektor.

\section*{$\mathbb{Z}_3$-Zirkulant}
$C = \bigl(\begin{smallmatrix}a&b&c\\c&a&b\\b&c&a\end{smallmatrix}\bigr)$,
Eigenwerte $\lambda_k = a + b\omega^k + c\omega^{2k}$, $\omega = e^{2\pi i/3}$.

\section*{$\tilde{T}\cdot m = 1$}
Kompaktifizierungsrelation: $R_m = \hbar/(mc^2)$.

\section*{$\xi$ aus Galois}
\[
\xi = \frac{(r_\mu/r_e)^2}{|GF(9)^*|^2 \cdot 5^2 \cdot |GF(27)|}
= \frac{144/25}{43200} = \frac{4}{30000}
\]

\chapter*{Anhang C — Wichtige Dokumente}
\addcontentsline{toc}{chapter}{Anhang C — Wichtige Dokumente}
\begin{center}
\small
\begin{tabular}{lp{7.5cm}}
\toprule
Dok. & Inhalt \\
\midrule
006 & Yukawa-Leiter, Wicklungszahlen $(r_i, p_i)$ \\
077 & $E=mc^2=E=m$, $c$ als Konvention \\
190 & Korrekturregister, offene Brücken \\
241 & Zwei-Schichten-Architektur \\
282 & $\mathbb{Z}_3$-Zirkulant, Massenoperator \\
307/334 & Zeit in QM, Superposition \\
312 & $\Lambda\to1/R_H^2$, $10^{123}$-Zerlegung \\
338 & $(m_\mu/m_e)^2=43200$ [K] \\
339 & Frobenius-Trennung massiv/masselos [B] \\
340 & $\Delta m^2$ aus $GF(27)^*$ [K] \\
342/343 & Primzahl-Harmonik, Orbifold-Zeta [B] \\
346--348 & Ladung, Gell-Mann-Nishijima, Generationen [B] \\
351/352 & Leptonreste, PDG-Modellabhängigkeit [K] \\
\bottomrule
\end{tabular}
\end{center}

\chapter*{Anhang D — Offene Brücken}
\addcontentsline{toc}{chapter}{Anhang D — Offene Brücken}
\begin{enumerate}
  \item \textbf{R113} [S]: Warum fallen $\varepsilon_i$ mit der Generation?
  \item \textbf{$\tau_\mu$ direkt} [S]: Konventionsfreier Test des schwachen Sektors.
  \item \textbf{P20} [S]: Kosmischer Exponent 10 vorwärts herleiten.
  \item \textbf{CMB-Peaks} [S]: Folge $\{1,6,14,26\}$ aus Torus-Moden.
  \item \textbf{Quarksektor} [S]: Koeffizienten rein geometrisch.
  \item \textbf{CKM-Phase} [S]: $\delta_{CP}\approx68°$ aus $GF(27)^*$.
\end{enumerate}


\chapter*{Anhang E — Teilchendiagramm und Belegstatus}
\addcontentsline{toc}{chapter}{Anhang E — Teilchendiagramm und Belegstatus (Dok. 356/357)}

\section*{Das Teilchendiagramm (Dok. 356)}

Dok.~356 ist ein Querformat-Übersichtsblatt das alle Standardmodell-Fermionen
und Eichbosonen als konzentrische Ring-Diagramme zeigt --- geordnet nach
Generationen und klassifiziert nach ihrer Galois-Herkunft.

\textbf{Aufbau jedes Teilchen-Diagramms} (von innen nach außen):
\begin{itemize}
  \item Weißer Innenkreis: Symbol und Ladung $Q$; links $Y_L$, rechts $Y_R$
  \item Farbring (nur Quarks): $r/g/b$ = Farbladung \status{[B] Dok.~346}
  \item Mittelring: violett links = $S_d$ = linkshändig; orange rechts = $S_u$ = rechtshändig \status{[B] Dok.~349}
  \item Außenring: Galois-Orbit-Farbe (blau O1/O3, hellblau O3, petrol O2, beere O4) \status{[B] Dok.~348}
  \item Rote Umrandung = $SU(2)_L$-Dublett-Paar; schwarze Umrandung = Generation 3
\end{itemize}

Das vollständige Diagramm ist unter
\url{https://github.com/jpascher/T0-Time-Mass-Duality} →
\texttt{2/pdf/356\_SM\_Teilchen\_GF27\_De.pdf} verfügbar.

\section*{Belegstatus der Quantenzahlen (Dok. 357)}

\subsection*{Algebraisch bewiesen [B]}

\begin{description}

\item[Elektrische Ladung $Q$] Dok.~346/347\\
$Q \in \{0, -1, +\tfrac{2}{3}, -\tfrac{1}{3}\}$ aus QR-Bit und $N \bmod 3$.
Alle vier Werte ohne freie Parameter.

\item[Hyperladung $Y$] Dok.~347 Satz~A\\
$Y = -1 + \tfrac{4}{3}N_\text{QR}$; $N_\text{QR}$ = Legendre-Symbol in $GF(27)^*$.

\item[Schwacher Isospin $I_3$] Dok.~347 Satz~B\\
$I_3 = \pm\tfrac{1}{2}$ aus $jE_7$-Projektor (Su/Sd-Trennung); Gradparität Cl$(6)$.

\item[Gell-Mann-Nishijima $Q = I_3 + Y/2$] Dok.~347 Satz~C\\
Vollständig algebraisch, kein freier Parameter.

\item[Farbladung $r/g/b$] Dok.~346 Satz~C\\
$Z_3$-Ladungsoperator $N \bmod 3$ in $GF(3)$. Drei Farben = drei Elemente.

\item[Chiralität $S_u / S_d$] Dok.~349\\
$S_u$ = gerade Cl$(6)$-Grade = rechtshändig;
$S_d$ = ungerade = linkshändig ($\gamma_5$-Eigenwert).

\item[$SU(2)_L$-Dublett] Dok.~349 Satz~D\\
Dublett = $S_d$-Sektor. Rechtshändiges Neutrino $\nu_R$ strukturell absent.

\item[3 Generationen] Dok.~348 Satz~A\\
Drei Frobenius-Zyklen in $GF(27)^*$. Ordnung 3 algebraisch erzwungen.

\item[Permutationsmatrix Quarks $\times$ Leptonen] Dok.~348 Satz~B\\
Spur-Bilinearform $GF(27)/GF(3)$: $f_3 \times f_4$-Struktur algebraisch bestimmt.

\item[Kopplungsgewichte $1/\sqrt{2}$] Dok.~348 Satz~C + Dok.~353\\
$f_1 \times f_3$-Bilinearform; zwei von drei Kopplungen aktiv.

\item[Koide-Amplitude $d/c = \sqrt{2}$] Dok.~353\\
$Z_3$-Gleichverteilung erzwingt $d/c = \sqrt{2}$ äquivalent zur Koide-Bedingung
$Q = 2/3$. Abgeschlossen 7.~Sept.~2026.

\item[8 Gluonen] Dok.~339\\
$4 \times Z_{13}$-Orbits $\times$ $2$ $Z_2$-Paritäten = 8. Strukturell erzwungen.

\end{description}

\subsection*{Numerisch bestätigt [K]}

\begin{description}
\item[Elektroschwache Gruppe $Z_2$] Dok.~336\\
$\mathrm{Gal}(\mathrm{GF}(9)/\mathrm{GF}(3)) = Z_2$. Numerisch bestätigt.
\end{description}

\subsection*{Offen [S]}

\begin{description}
\item[Generations-Zuordnung (Reihenfolge)] Dok.~346\\
Welches Orbit-Element = Generation 1/2/3? Orbits algebraisch gleichwertig;
Zuordnung nicht abgeleitet.

\item[CKM/PMNS-Winkel vollständig] Dok.~348 Satz~D\\
Struktur bekannt [B]; vollständige Winkelwerte erfordern weitere Arbeit.
\end{description}

\subsection*{Experimentell gesichert, FFGFT-Herleitung offen [X]}

\begin{description}
\item[Higgs $H^0$] Dok.~019\\
$\tilde{T} \cdot m = 1$ ersetzt den Higgs-Mechanismus konzeptuell.
Experimentell gesichert (CERN 2012); formale FFGFT-Herleitung offen.
\end{description}

\end{document}
